\newpage
\section{Conclusão}

O trabalho reuniu os principais conteúdos estudados na disciplina em uma aplicação de pintura robotizada. A partir da escolha do Stäubli TX90, foram desenvolvidas a modelagem cinemática, a geração da trajetória, a verificação do espaço de trabalho, o cálculo das velocidades, a dinâmica inversa e o controle do movimento.

A trajetória final foi concluída em $67{,}80\,\text{s}$ e todas as poses apresentaram solução de cinemática inversa. Os limites de velocidade e torque foram respeitados, sendo o maior esforço de $105{,}50\,\text{Nm}$ em J2. Entre os controladores avaliados, o torque computado apresentou os menores erros RMS, embora o PD com compensação de gravidade também tenha acompanhado a referência com erros reduzidos.

Dessa forma, os resultados indicam que a tarefa é viável no modelo utilizado. Para uma implementação física, ainda seriam necessários testes de colisão, calibração dos referenciais, identificação da ferramenta real e análise do processo de deposição da tinta. Mesmo com essas limitações, o modelo desenvolvido forma uma base que pode ser adaptada para o preenchimento da bandeira e para outras tarefas dentro da mesma linha de produção.
